% talk.tex — the deck, rendered through the Cyberwitchery beamer theme.
% Slide body lives in content.tex.
%
% Build:  make     (or: xelatex talk)
\documentclass[aspectratio=169,14pt]{beamer}

\usetheme{Cyberwitchery}        % dark; add [light] for the light palette

% ------------------------------------------------------------------
% Map the old deck's accent-colour names onto the Cyberwitchery palette,
% so content.tex (which references copper/steel/mute/ink/paper/rule) renders
% unchanged.
% ------------------------------------------------------------------
\colorlet{copper}{cwTeal}
\colorlet{steel}{cwTeal}
\colorlet{mute}{cwFaint}
\colorlet{ink}{cwInk}
\colorlet{paper}{cwGround}
\colorlet{rule}{cwRule}

% ------------------------------------------------------------------
% Readability: ordinary explanatory bullets read as prose, not mono, so the
% mono voice is reserved for titles, labels, code, short claims, and emphasis.
% (The theme already scales \cwhanken to the mono and routes its bold to mono
% bold; \cwhankenbody is just a name this deck's content still references.)
% ------------------------------------------------------------------
\newcommand{\cwhankenbody}{\cwhanken}
\setbeamertemplate{itemize/enumerate body begin}{\cwhankenbody}
\setbeamertemplate{itemize/enumerate subbody begin}{\cwhankenbody}
\setbeamertemplate{itemize/enumerate subsubbody begin}{\cwhankenbody}

% ------------------------------------------------------------------
% The five content macros, reimplemented in the theme's idiom.
% (Same names/signatures as talk.tex, so the shared body is verbatim.)
% ------------------------------------------------------------------
% Emphasis: bold + teal in the surrounding face. In a Hanken bullet that bold is
% routed to JetBrains Mono Bold (via cwhankenbody), so it stays mono and sized to
% the body; in a mono hero it is full-size mono bold. No fixed \sffamily, which
% would force full-size mono and tower over scaled body text.
\newcommand{\copperaccent}[1]{\textcolor{cwTeal}{\textbf{#1}}}

\newcommand{\hero}[1]{%
  \vfill
  \begin{center}{\bfseries\large\color{cwInk} #1\par}\end{center}
  \vfill}

\newcommand{\subhero}[1]{%
  \begin{center}{\cwhankenbody\color{cwMuted} #1\par}\end{center}}

% A tour stop: code on top, a small note of which notation it is, then the
% claim and the attribution.  #1 claim  #2 code  #3 notation  #4 attribution
\newcommand{\claim}[4]{%
  \vfill
  \begin{center}
    {\ttfamily\large\color{cwTeal} #2\par}
    \vspace{0.35em}
    {\cwhankenbody\footnotesize\color{cwMuted} #3\par}
    \vspace{1.0em}
    {\bfseries\Large\color{cwInk} #1\par}
    \vspace{1.4em}
    {\cwhankenbody\footnotesize\color{cwFaint} #4\par}
  \end{center}
  \vfill}

% Full-bleed live-coding cue.
\newcommand{\livecode}[1]{%
  \begin{frame}[plain, noframenumbering]
    \begin{tikzpicture}[remember picture, overlay]
      \node[anchor=west, text=cwTeal, font=\bfseries\large]
        at ([xshift=0.06\paperwidth, yshift=2.6cm]current page.west)
        {$\rightarrow$\ live};
      \node[anchor=west, text=cwInk, font=\bfseries\fontsize{40}{48}\selectfont,
            text width=0.82\paperwidth]
        at ([xshift=0.06\paperwidth, yshift=-0.2cm]current page.west)
        {#1};
    \end{tikzpicture}
  \end{frame}}

% ------------------------------------------------------------------
\title{type systems\\ you might not know}
\subtitle{(but will love)}
\author{veit heller}
\institute{cyberwitchery lab}
\date{2026}

\begin{document}

  % --- title : faint logo on the right, no footline ---
  {\setbeamertemplate{footline}{}
  \begin{frame}[plain]
    \cwlogoright[cwGhost]{0.42\paperheight}%
    \titlepage
  \end{frame}}

    % ==================================================================
  \section{beyond static vs dynamic}
  % ~ 4 min --- earn the thesis, don't just state it
  % ==================================================================
  \begin{frame}{the menu you were handed}
    \begin{itemize}
      \item static vs dynamic.
      \item strong vs weak.
      \item \textcolor{mute}{most of us stopped there.}
    \end{itemize}
  \end{frame}

  \begin{frame}{the usual two axes}
    \hero{static vs dynamic is about \emph{when} checks run.\\[0.3em]
    strong vs weak barely has a consistent meaning.}
    \subhero{neither is the interesting axis.}
  \end{frame}

  \begin{frame}{a better question}
    \hero{i think it's more useful to ask\\[0.3em]
    \copperaccent{what a type can say} about your program.}
  \end{frame}

  \begin{frame}{what can a type say?}
    \begin{itemize}
      \item[\textcolor{mute}{\textbar}] \textcolor{mute}{“this is an integer.”}
      \item[\textcolor{mute}{\textbar}] \textcolor{mute}{“this is a list of strings.”}
      \item “this integer is never zero.”
      \item “this list is non-empty, so \texttt{head} is safe.”
      \item “this socket has already been closed.”
      \item “this function touches the disk.”
      \item “send a request, then receive exactly one reply.”
    \end{itemize}
    \vspace{0.4em}
    \subhero{the greyed-out two, you say every day. the rest is today's talk.}
  \end{frame}

  \begin{frame}{the plan}
    \begin{itemize}
      \item five things a type can say, \copperaccent{from familiar to strange}.
      \item the first is probably familiar. the last will not be.
      \item a roadmap, so you leave with somewhere to go.
    \end{itemize}
    \subhero{one trip to the terminal, to build a type checker ourselves.}
  \end{frame}

  % ==================================================================
  \section{five things a type can say}
  % ~ 14 min --- one ladder, five rungs; rung 2 (refinement) is live
  % ==================================================================
  % --- rung 1: ownership in rust, the warmup ------------------------
  \begin{frame}{ownership}
    \hero{“this value has \copperaccent{exactly one owner}.”}
    \subhero{use it once, then move it or borrow it. rust's borrow checker is
    one of these, whether or not you have met it.}
  \end{frame}

  \begin{frame}[fragile]{can't close it twice}
    \vfill
\begin{lstlisting}[morekeywords={let}]
let c = Conn::open("socket-7");
c.close();
c.close();
\end{lstlisting}
    \vskip6pt
    {\ttfamily\small\textcolor{copper}{error[E0382]: use of moved value: `c`}}\\[2pt]
    {\ttfamily\small\textcolor{mute}{note: `close` takes ownership of the receiver `self`, which moves `c`}}
    \vfill
    \subhero{close consumes the handle, so a second use can't typecheck.}
  \end{frame}

  \begin{frame}{that's an “affine” type}
    \begin{itemize}
      \item \textbf{linear}: use a value \emph{exactly} once.
      \item \textbf{affine}: use it \emph{at most} once.
      \item \textbf{rust}: affine, unless a type is \texttt{Copy} (values safe to duplicate).
    \end{itemize}
    \subhero{a large class of use-after-free, double-free, and data-race bugs
    becomes a compile error. \texttt{Send} and \texttt{Sync} carry thread-safety in the type.}
  \end{frame}

  % --- rung 2: refinement, the live build --------------------------
  \livecode{let's write a type checker.}

  \begin{frame}[t]{a real type system}
    \vskip0.09\paperheight
    \cwstat{55}{lines, no solver}
    \vskip0.07\paperheight
    {\color{cwRule}\hrule height0.6pt}\vskip10pt
    {\cwhankenbody\color{cwFaint}\small that is a working refinement type checker, in plain interval arithmetic. real systems (liquid haskell, f\textsuperscript{*}) hand the subtyping to an smt solver.}
  \end{frame}

  \begin{frame}{effect systems}
    \claim{“this function touches the disk.”}
      {readFile : Path ->\{IO\} Text}
      {unison}
      {in koka, eff, and unison, the effect rides in the type next to the return value.}
  \end{frame}

  \begin{frame}{session types}
    \claim{“the worker can't get a job before it registers.”}
      {!Register ; ?Job ; !Result ; Close}
      {freest}
      {two concurrent processes talking over one channel. in freest the order is part of the channel's type, so sending them out of order won't compile.}
  \end{frame}

  \begin{frame}{dependent types}
    \claim{“this list is non-empty, so \texttt{head} is always safe.”}
      {head : Vec A (suc n) $\to$ A}
      {agda}
      {types can depend on \emph{values}. \texttt{Vec A n} carries its length,
       so \texttt{head []} does not typecheck.}
  \end{frame}

  \begin{frame}{proofs as programs}
    \claim{“concatenation adds the lengths.”}
      {\_++\_ : Vec A m $\to$ Vec A n $\to$ Vec A (m + n)}
      {agda}
      {the signature is the spec, and the two-line body is a proof of it,
       checked at compile time.}
  \end{frame}

  % --- synthesis: the five rungs, one move -------------------------
  \begin{frame}[t]{five claims, side by side}
    \vskip8pt
    \footnotesize
    \renewcommand{\arraystretch}{1.4}%
    \arrayrulecolor{cwRule}\setlength{\arrayrulewidth}{0.6pt}%
    \noindent\begin{tabular}{@{}>{\raggedright\arraybackslash}m{0.17\textwidth}
                              >{\raggedright\arraybackslash}m{0.33\textwidth}
                              >{\columncolor{cwPanel}\raggedright\arraybackslash}m{0.36\textwidth}@{}}
      {\cwhankenbody\color{cwMuted}} & {\cwhankenbody\color{cwMuted}the bug, at runtime} & {\cwhankenbody\color{cwInk}\bfseries caught at compile time} \\
      \hline
      {\cwhankenbody\color{cwFaint}ownership}  & {\cwhankenbody use after close, double free} & {\cwhankenbody moved, so reuse won't compile} \\
      \hline
      {\cwhankenbody\color{cwFaint}refinement} & {\cwhankenbody divide by zero} & {\cwhankenbody argument must be non-zero} \\
      \hline
      {\cwhankenbody\color{cwFaint}effects}    & {\cwhankenbody a hidden disk write} & {\cwhankenbody the effect is in the signature} \\
      \hline
      {\cwhankenbody\color{cwFaint}session}    & {\cwhankenbody a message out of order} & {\cwhankenbody wrong message won't compile} \\
      \hline
      {\cwhankenbody\color{cwFaint}dependent}  & {\cwhankenbody \texttt{head []}, off by one} & {\cwhankenbody the length lives in the type} \\
    \end{tabular}
  \end{frame}

  % ==================================================================
  \section{a roadmap}
  % ~ 3 min
  % ==================================================================
  \begin{frame}{pick an axis, not a language}
    \begin{itemize}
      \item resource safety $\to$ rust.
      \item proof $\to$ agda, idris, lean.
      \item effects $\to$ koka, unison.
      \item contracts $\to$ liquid haskell, f\textsuperscript{*}.
    \end{itemize}
    \subhero{you don't need a phd to use any of these in anger.}
  \end{frame}

  \begin{frame}{resources}
    \begin{itemize}
      \item \textit{types and programming languages}, pierce.
      \item \textit{programming language foundations in agda}.
      \item \textit{the little typer}, friedman \& christiansen.
    \end{itemize}
    \subhero{none of them will melt your brain.}
  \end{frame}

  \begin{frame}{the takeaway}
    \hero{a type system is a language\\
    for saying what's true.\\[0.4em]
    pick one that lets you say\\
    \copperaccent{what you mean}.}
  \end{frame}


  % --- closing ---
  {\setbeamertemplate{footline}{}
  \begin{frame}[plain]
    \cwlogoright[cwInk]{0.42\paperheight}%
    \vskip0.22\paperheight
    \begin{minipage}{0.62\textwidth}
      {\usebeamerfont{title}\color{cwInk}thanks.}\par\vskip12pt
      {\cwhanken\large\color{cwTeal}questions?}\par
      \vskip0.10\paperheight
      {\cwhanken\color{cwMuted}veit heller \quad \texttt{veit@veitheller.de}}\\[5pt]
      {\cwhanken\color{cwFaint}slides: \texttt{github.com/hellerve/talks}}
    \end{minipage}
  \end{frame}}

\end{document}
